%&latex
% !TeX program = lualatex
% Use A4 explicitly so output is consistent across systems.
\documentclass[a4paper,12pt]{article}

% Page layout.
\usepackage[margin=1in]{geometry}
\setlength{\parindent}{2em}
\setlength{\parskip}{0pt}

% Fonts.
\usepackage{fontspec}
\usepackage[fontset=fandol]{ctex}
\setmainfont{TeX Gyre Termes}
\setmonofont{TeX Gyre Cursor}

% Syntax-highlighted code blocks.
% Compile with shell escape enabled when using minted.
\usepackage{xcolor}
\usepackage{minted}
\setminted{
  fontsize=\small,
  breaklines=true,
  autogobble=true,
  frame=single,
  framesep=2mm
}

% Clickable cross-references and links. Keep this near the end of the preamble.
\usepackage[colorlinks=true, linkcolor=blue, urlcolor=blue]{hyperref}

\title{工作坊讲义：本地运行大语言模型}
\author{Yang Shen}
\date{}

\begin{document}

\maketitle

\section{概述}

本文档演示如何在本地运行小型语言模型，并通过智能体命令行工具使用它。示例中使用 \texttt{llama.cpp} 作为本地推理运行时，Qwen3-0.6B 作为练习用的小型模型，以及 \texttt{pi} 作为智能体式交互的界面。

本教程的目标不是构建速度最快或能力最强的本地 AI 系统，而是学习通用工作流：安装所需的工具、下载模型文件、启动本地模型服务器、将应用连接到它，并尝试一些任务。模型被故意选得较小，以便在普通个人电脑上即可操作。

本文档分为两个主要部分。快速入门部分给出从设置到运行本地模型的最短路径。后面的章节更详细地解释工具、文件格式、命令以及各种权衡。

\end{document}
